\documentclass[12pt,letterpaper]{article}

\usepackage[utf8]{inputenc}
\usepackage[spanish]{babel}
\usepackage[T1]{fontenc}
\usepackage{amsmath}
\usepackage{csquotes} %

\usepackage[
    top=3cm,
    left=4cm,
    bottom=3cm,
    right=2cm
]{geometry}

\usepackage{setspace}
\onehalfspacing

\usepackage{mathptmx} 

\usepackage{fancyhdr}
\pagestyle{fancy}
\fancyhf{}
\fancyfoot{}
\fancyhead[R]{\thepage}
\renewcommand{\headrulewidth}{0pt}

\usepackage{graphicx}
\usepackage{caption}
\captionsetup{font=small, labelfont=bf}

\usepackage[backend=biber, style=ieee, sorting=none]{biblatex}
\addbibresource{references.bib}   

\usepackage{titlesec}
\titleformat{\section}{\normalfont\bfseries\large}{\thesection.}{1em}{\MakeUppercase}
\titleformat{\subsection}{\normalfont\bfseries}{\thesubsection}{1em}{\MakeUppercase}
\titleformat{\subsubsection}{\normalfont\bfseries\normalsize}{\thesubsubsection}{1em}{\MakeUppercase}

\begin{document}
\begin{titlepage}
    \thispagestyle{empty}
    \centering
    \vspace*{2cm}
    {\LARGE \textbf{CHATBOTS CONVERSACIONALES VS. FORMULARIOS TRADICIONALES EN LA UPTC}}\\[3cm]
    {\large Edwin David Fino Velandia\\ Edward Mauricio Guatibonza Suárez\\ Fredy Alejandro Bonilla Rodríguez}\\[2cm]
    {\large Universidad Pedagógica y Tecnológica de Colombia (UPTC)}\\
    {\large Facultad de Ingeniería}\\
    {\large Escuela de Ingeniería de Sistemas y Computación}\\[2cm]
    {\large Tunja}\\
    {\large \the\year}
\end{titlepage}

\newpage

\thispagestyle{empty}
\tableofcontents
\newpage

\setcounter{page}{1}

Texto aquí...

\section{Planteamiento del Problema}

\subsection{Antecedentes}

En las instituciones de educación superior, la recolección de información ha pasado progresivamente de formatos físicos a herramientas digitales \cite{kim2019, hernandez2014}. A pesar de este avance, los formularios web tradicionales todavía presentan dificultades importantes, como bajos niveles de respuesta y abandono antes de completar las preguntas. Por esta razón, se ha empezado a estudiar el uso de chatbots en aplicaciones de mensajería instantánea como una alternativa que podría facilitar la participación de los estudiantes \cite{xiao2020, beam2023}.

Xiao et al. \cite{xiao2020} desarrollaron un estudio de campo con cerca de 600 participantes con el fin de evaluar la recolección de datos en encuestas conversacionales administradas por un chatbot impulsado por inteligencia artificial frente a un formulario web tradicional. Los autores hallaron que la tasa de finalización obtenida en la modalidad de chatbot conversacional alcanzó el 54\%, superando a más del doble la registrada en el formulario web estático, la cual fue del 24.2\%. Este antecedente demuestra empíricamente que la interacción guiada pregunta a pregunta incrementa el compromiso del participante \cite{xiao2020}.

Por su parte, Beam \cite{beam2023} evaluó experimentalmente la efectividad de las redes sociales y agentes conversacionales como herramientas de reclutamiento y recolección de información. Mediante un diseño experimental aleatorizado aplicado a encuestas docentes, el estudio demostró que los chatbots alcanzaron tasas de respuesta y una calidad de datos significativamente superiores en comparación con los formularios web estáticos convencionales.

Sin embargo, Zarouali et al. \cite{zarouali2023} llevaron a cabo una investigación longitudinal comparando el comportamiento de respuesta, la calidad de datos y la evaluación de usuario entre chatbots y encuestas web. A diferencia de los hallazgos de Xiao et al. y Beam, esta investigación no encontró ventajas estadísticamente significativas en el uso del chatbot frente al formulario tradicional en el tiempo, e incluso registró una disminución en la valoración favorable por parte de los usuarios en mediciones repetidas. Este resultado evidencia la falta de consenso en la comunidad científica y la necesidad de evitar asumir resultados de antemano sin una contrastación local.

Kim et al. \cite{kim2019} analizaron el impacto de la plataforma y el estilo conversacional sobre la calidad de las respuestas en encuestas en línea. Los autores concluyeron que los formularios web tradicionales operan como sistemas receptores pasivos que no confirman en tiempo real la recepción de datos. Esta estructura estática entra en conflicto directo con las expectativas de interactividad inmediata que poseen los usuarios habituados al entorno digital de la mensajería instantánea.

En cuanto a la extensión del instrumento, Hoerger \cite{hoerger2010} examinó la deserción de participantes en función de la longitud del cuestionario en entornos universitarios. La evidencia cuantitativa demonstrated que cerca del 10\% de los estudiantes abandona la prueba de inmediato al percibir una interfaz extensa, y que la tasa de deserción se incrementa aproximadamente 2 puntos porcentuales por cada 100 preguntas adicionales.

Desde la perspectiva del diseño de interfaz, Nielsen y Molich \cite{nielsen1990} establecieron los principios de la evaluación heurística. Sus postulados explican cómo la falta de retroalimentación inmediata, la escasa visibilidad del estado del sistema y la ausencia de indicadores de progreso en formularios estáticos generan carga cognitiva e inducen al abandono prematuro.

Asimismo, Harley et al. \cite{harley2007} evaluaron el uso de la mensajería móvil en la comunicación académica universitaria. Los resultados revelaron que la comunidad estudiantil otorga prioridad y revisión continua a las aplicaciones de mensajería instantánea por encima de los canales formales como el correo electrónico, donde las notificaciones tienden a postergarse e ignorarse.

En el diseño de la interacción conversacional, Araujo \cite{araujo2018} demostró que las señales de diseño antropomórfico y el estilo comunicativo de un agente conversacional influyen directamente en la percepción de calidez y competencia que el usuario atribuye tanto al chatbot como a la organización que lo despliega. Este hallazgo refuerza la idea de que no basta con automatizar la conversación, sino que su estilo debe diseñarse cuidadosamente para maximizar el compromiso del respondiente.

En un entorno clínico-administrativo, te Pas et al. \cite{tepas2020} compararon la experiencia de usuario de un cuestionario tipo chatbot frente a un cuestionario computarizado tradicional. Los resultados mostraron una valoración mixta: el 46.7\% de los términos evaluados favorecieron al chatbot, el 8\% favorecieron al formulario tradicional y el 48\% restante no mostró diferencias significativas, sin que se registraran variaciones relevantes en el tiempo de finalización entre ambas modalidades.

De manera más directamente comparable al presente estudio, Çavuşoğlu Deveci et al. \cite{devecifuchs2026} realizaron un experimento de campo aleatorizado con estudiantes universitarios, comparando una interfaz tipo chatbot frente a un formulario web tradicional. Los resultados indicaron que los estudiantes percibieron el diseño tipo chatbot como más original y entretenido, aunque también como más difícil de navegar; no se hallaron diferencias significativas en el tiempo de respuesta ni en la extensión de las respuestas abiertas, mientras que la proporción de encuestados con datos faltantes fue marginalmente menor en el chatbot.

En el ámbito metodológico, Hernández-Sampieri et al. \cite{hernandez2014} advierten que una alta tasa de deserción no aleatoria en la recolección de datos genera un sesgo de selección. Este fenómeno compromete la validez externa de los resultados e impide la representatividad de la muestra respecto a la población de estudio.

En el contexto colombiano, Barón Salazar et al. \cite{baronsalazar2026} documentaron cómo la dependencia de procesos manuales en oficinas universitarias genera cuellos de botella que incrementan el riesgo de deserción por falta de respuesta inmediata. En su estudio, la automatización del canal conversacional con Procesamiento de Lenguaje Natural redujo el tiempo de respuesta en un 99.81\% frente a la atención presencial.

Por otro lado, Cisternas-Osorio et al. \cite{cisternas2022} analizaron el uso de Telegram como canal de comunicación institucional en universidades hispanohablantes. Los autores encontraron que en Latinoamérica el uso de Telegram en el ámbito universitario es todavía minoritario y que los canales existentes operan de manera predominantemente unidireccional. Este hallazgo respalda la pertinencia de explorar aplicaciones interactivas de Telegram en universidades colombianas, dado que la plataforma permanece subutilizada en este contexto.

Esta tensión teórica fundamenta la necesidad de ejecutar un experimento comparativo en el contexto de la Universidad Pedagógica y Tecnológica de Colombia (UPTC). El presente estudio busca evaluar si la implementación de un chatbot conversacional en Telegram logra mitigar la brecha de canal y la fricción de usabilidad que afectan a los formularios web institucionales \cite{baronsalazar2026, cisternas2022}.

\subsection{Descripción de la situación problemática}
Para el periodo 2025--2026, la UPTC depende únicamente del correo institucional y de formularios web estáticos para realizar sus encuestas. Debido a que la mayor parte de los estudiantes utiliza habitualmente servicios de mensajería instantánea, el envío de formularios extensos es ignorado por un amplio sector estudiantil o abandonado a mitad del proceso, al ser percibido como un trámite rígido u obligatorio más que como un espacio de participación.

\subsection{Formulación del Problema}

Con base en el diagnóstico y la situación problemática identificada en la UPTC, se plantea la siguiente pregunta central de investigación:

\begin{quote}
¿En qué medida la implementación de un chatbot conversacional e interactivo en Telegram influye en la tasa de finalización, el tiempo de respuesta y la usabilidad percibida en la recolección de datos estudiantiles en la UPTC, en comparación con los formularios web estáticos tradicionales?
\end{quote}

Para abordar de manera integral esta interrogante central, el problema se desglosa en las siguientes preguntas específicas:

\begin{enumerate}
    \item ¿Cómo debe diseñarse y desarrollarse un prototipo de chatbot conversacional en Telegram, integrado con Python, para la administración efectiva de instrumentos de recolección de datos en la UPTC?
    \item ¿Existen diferencias estadísticamente significativas en la tasa de finalización, la tasa de abandono y el tiempo de respuesta entre la recolección mediante el chatbot conversacional y el formulario web tradicional?
    \item ¿Cuál es la evaluación de los estudiantes sobre la usabilidad percibida, la carga cognitiva y la experiencia de usuario al interactuar con la interfaz conversacional en comparación con la interfaz web estática?
\end{enumerate}

\subsection{Consecuencias de no Investigar}
De mantenerse esta tendencia durante el periodo 2026--2027, el abstencionismo estudiantil en los procesos de consulta de la UPTC se profundizará. Como consecuencia, las decisiones administrativas, académicas y de bienestar institucional se seguirán formulando a partir de muestras reducidas e irrepresentativas (con tasas de participación inferiores al 23\% registrado a nivel global), perpetuando el sesgo en la evaluación docente y distorsionando la asignación de recursos hacia proyectos que no reflejan las necesidades reales de la comunidad universitaria.

\subsection{Delimitación}
Para mitigar esta problemática, se propone el diseño y evaluación experimental de un chatbot conversacional e interactivo en Telegram (desarrollado en Python) para la aplicación de encuestas institucionales. El objetivo es determinar si la integración de una plataforma de uso cotidiano reduce la fricción de interfaz, incrementa la tasa de respuesta y finalización, y disminuye el tiempo de diligenciamiento en comparación con los formularios web tradicionales.

\section{Justificación}

La realización de este proyecto de investigación se fundamenta en la necesidad de abordar los bajos índices de participación estudiantil en la recolección de información institucional dentro de la Universidad Pedagógica y Tecnológica de Colombia (UPTC). A continuación, se detallan las razones teóricas, metodológicas, prácticas y tecnológicas que justifican la ejecución de este estudio:
\paragraph{Justificación Teórica}
En la literatura especializada existe una controversia activa sobre la efectividad de los agentes conversacionales frente a los formularios estáticos. Mientras que diversas investigaciones afirman que las interfaces conversacionales incrementan significativamente las tasas de respuesta, otros estudios reportan la ausencia de ventajas estadísticamente significativas o incluso un incremento en la dificultad de navegación. Esta investigación aporta evidencia empírica cuantitativa obtenida directamente en el contexto universitario colombiano, enriqueciendo el cuerpo de conocimiento sobre interacción humano-computador y comunicación digital institucional en América Latina.

\paragraph{Justificación Metodológica}
El proyecto propone un diseño experimental riguroso para comparar objetivamente dos canales de interacción digital bajo condiciones reales de uso. Mediante la medición estandarizada de métricas objetivas de rendimiento (tiempo de diligenciamiento y tasa de finalización) y la evaluación de la usabilidad percibida (mediante escalas validadas como SUS o UEQ), se establece un marco evaluativo replicable para la validación de interfaces conversacionales en el área de la ingeniería de software.

\paragraph{Justificación Práctica e Institucional}
A nivel institucional, la UPTC requiere superar la dependencia de canales pasivos y poco efectivos que generan sesgos de selección en las evaluaciones docentes, estudios de bienestar y diagnósticos académicos. La implementación de un mecanismo interactivo ajustado a los hábitos digitales cotidianos de la población estudiantil fortalece los canales de escucha activa, fomenta la participación universitaria y proporciona a las directivas datos más representativos para la toma de decisiones.

\paragraph{Justificación Tecnológica}
Desde la perspectiva de la Ingeniería de Sistemas, este trabajo demuestra la viabilidad técnica y operativa de integrar la API de Telegram con entornos de desarrollo en Python para construir asistentes de captura de datos ligeros, multiplataforma, de bajo costo y fácil mantenimiento, reduciendo la fricción asociada al uso de aplicaciones propietarias o formularios web extensos.

\section{Objetivos}
\subsection{Objetivo General}
Contribuir a que la UPTC cuente con criterios para decidir si incorpora canales conversacionales en sus procesos institucionales de consulta, de manera que se aumente la participación estudiantil y la representatividad de la información con la que toma decisiones.

\subsection{Objetivos Específicos}
\begin{enumerate}
  \item Caracterizar el perfil de los estudiantes de los dos grupos para valorar su comparabilidad. \textit{Meta:reportar el perfil de comparabilidad de una muestra de aproximadamente 60 estudiantes (30 por grupo).}
  \item Determinar la tasa de respuesta obtenida con el chatbot de Telegram y con el formulario web. \textit{Meta: reportar el porcentaje de estudiantes invitados que respondió, por grupo.}reportar el porcentaje de estudiantes invitados que respondió, por grupo.
  \item Determinar la tasa de finalización y el punto de abandono en ambas modalidades. \textit{}Meta:reportar el porcentaje de encuestas completadas frente a las iniciadas y la pregunta con mayor abandono, por grupo.
  \item Comparar las tasas de respuesta y de finalización entre las dos modalidades para estimar el efecto del canal. \textit{Meta: reportar la diferencia entre grupos y su incertidumbre para ambas tasas.}
\end{enumerate}

\section{Alcance de la investigación}

La presente investigación se clasifica dentro de un alcance explicativo y correlacional, con un componente descriptivo previo, de acuerdo con la tipología de Hernández Sampieri (Capítulo 5):

\begin{enumerate}
\item \textbf{Alcance Correlacional:} El estudio evalúa el grado de asociación y la relación directa entre la variable independiente (Canal de recolección de información: Chatbot interactivo en Telegram vs. Formulario web estático) y las variables dependientes (Tasa de finalización de la encuesta, Tiempo promedio de respuesta y Percepción de usabilidad del usuario).

\item \textbf{Alcance Explicativo:} La investigación va más allá de la caracterización del fenómeno, ya que busca determinar relaciones de causa-efecto mediante un diseño quasi-experimental. Se pretende verificar si la reducción de la fricción de interfaz al migrar las consultas a un entorno conversacional cotidiano es la causa directa del incremento significativo en la tasa de respuesta y de la disminución del tiempo de diligenciamiento en los estudiantes.

\item \textbf{Componente Descriptivo:} Como fase previa al análisis relacional, se caracterizan y miden cuantitativamente las propiedades de interacción, el porcentaje de abandono y la distribución de tiempos de respuesta en cada uno de los dos medios evaluados.

\end{enumerate}


\section{Formulación de Hipótesis}
\textbf{H$_{1a}$:} los estudiantes que reciben la encuesta institucional mediante un chatbot de Telegram presentan una tasa de respuesta mayor que los estudiantes que la reciben mediante el formulario web estático.

\textbf{H$_{0a}$:} no existe diferencia en la tasa de respuesta entre los estudiantes que reciben la encuesta por chatbot de Telegram y por formulario web.

\textbf{H$_{1b}$:} los estudiantes que responden la encuesta mediante el chatbot de Telegram presentan una tasa de finalización mayor que los que la responden mediante el formulario web estático.

\textbf{H$_{0b}$:} no existe diferencia en la tasa de finalización entre los estudiantes que responden por chatbot de Telegram y por formulario web.

\section{Variables}

En el marco del diseño experimental planteado, las variables de estudio se clasifican según su función dentro de la investigación y se definen conceptual y operacionalmente a continuación:

\subsection{Variable Independiente (VI)}

\begin{itemize}
    \item \textbf{Tipo de interfaz de recolección de datos:} Medio o canal digital mediante el cual el participante interactúa para responder el cuestionario institucional. Esta variable se manipula en dos niveles de tratamiento:
    \begin{itemize}
        \item \textit{Tratamiento experimental:} Interfaz conversacional e interactiva administrada por un chatbot en Telegram (desarrollado en Python).
        \item \textit{Tratamiento de control:} Interfaz web estática tradicional (formulario web convencional).
    \end{itemize}
\end{itemize}

\subsection{Variables Dependientes (VD)}

\begin{itemize}
    \item \textbf{Tasa de finalización ($VD_1$):} Porcentaje de participantes que completan y envían la totalidad del cuestionario con respecto al total de personas que inician la interacción.
    \item \textbf{Tiempo de diligenciamiento ($VD_2$):} Duración total, medida en segundos, desde el instante en que el participante inicia la primera pregunta hasta la confirmación de envío final.
    \item \textbf{Usabilidad percibida ($VD_3$):} Grado de facilidad de uso, eficiencia y satisfacción del usuario al interactuar con el sistema, evaluado formalmente a través de la escala estandarizada \textit{System Usability Scale} (SUS).
    \item \textbf{Tasa de abandono o deserción ($VD_4$):} Porcentaje de participantes que abandonan el flujo de interacción antes de responder la última pregunta o que dejan ítems obligatorios incompletos.
\end{itemize}

\subsection{Variables de Control y Extrañas}

\begin{itemize}
    \item \textbf{Contenido y extensión del instrumento:} Se garantiza que el número de preguntas, la redacción, las opciones de respuesta y el nivel de complejidad sean idénticos en ambos tratamientos.
    \item \textbf{Perfil demográfico de los participantes:} Estudiantes matriculados en la Facultad de Ingeniería de la UPTC, asegurando relativa homogeneidad en competencias digitales previas.
    \item \textbf{Dispositivo de acceso:} Registro del tipo de hardware empleado (dispositivo móvil o computador de escritorio) para controlar y analizar posibles sesgos en la velocidad de respuesta.
\end{itemize}

\subsection{Matriz de Operacionalización de Variables}

\begin{table}[htbp]
\centering
\caption{Matriz de Operacionalización de Variables}
\label{tab:operacionalizacion}
\small
\begin{tabular}{|p{3cm}|p{3.5cm}|p{3.5cm}|p{2.5cm}|p{2.5cm}|}
\hline
\textbf{Variable} & \textbf{Definición Conceptual} & \textbf{Indicador} & \textbf{Escala de Medición} & \textbf{Instrumento / Fuente} \\ \hline
\textbf{Tipo de Interfaz (VI)} & Canal digital empleado para la captura de información. & Tipo de tratamiento asignado (Chatbot vs. Formulario Web). & Nominal / Categórica & Registro del sistema / Logs \\ \hline
\textbf{Tasa de Finalización ($VD_1$)} & Proporción de encuestas respondidas en su totalidad. & $\frac{\text{Encuestas completas}}{\text{Total iniciadas}} \times 100$ & Razón / Porcentaje (\%) & Logs del servidor / Base de datos \\ \hline
\textbf{Tiempo de Respuesta ($VD_2$)} & Tiempo invertido en diligenciar el cuestionario. & $t_{\text{final}} - t_{\text{inicio}}$ (segundos) & Razón / Continua & Registros de tiempo (\textit{timestamps}) \\ \hline
\textbf{Usabilidad Percibida ($VD_3$)} & Grado de facilidad y satisfacción de uso. & Puntaje global SUS (Escala de 0 a 100). & Intervalo & Cuestionario SUS post-prueba \\ \hline
\textbf{Tasa de Abandono ($VD_4$)} & Deserción prematura antes de completar la prueba. & $\frac{\text{Sesiones abandonadas}}{\text{Total iniciadas}} \times 100$ & Razón / Porcentaje (\%) & Logs de interacción \\ \hline
\end{tabular}
\end{table}

\section{Marco Referencial}

\subsection{Marco Teórico Conceptual}

\subsubsection{Pruebas de Software}
Texto aquí...

\subsubsection{Automatización de Pruebas}
Texto aquí...

\subsubsection{Priorización de Pruebas}
Texto aquí...

\subsection{Marco Histórico}
Texto aquí...

\subsection{Estado Actual del Tema}
Texto aquí...

\subsection{Marco Científico y Tecnológico}

\subsubsection{Fundamentos Científicos}
Texto aquí...

\subsubsection{Tecnologías Relevantes}
Texto aquí...

\subsection{Marco Legal y Normativo}
Texto aquí... 

\section{Diseño Metodológico}
Texto aquí...

\section{Resultados y Discusión}
Texto aquí...

\section{Conclusiones y Recomendaciones}
Texto aquí...

\printbibliography[title={Bibliografía}]

\section*{Anexos}
Texto aquí...

\end{document}
